\subsection{RQ2: Non-IID Robustness.}
\label{sec:rq2}

To answer RQ2, we evaluate P4P in heterogeneous, non-IID environments. The objective is twofold: (1) to verify its defense efficacy under severe data distribution skew, and (2) to ensure it does not falsely penalize benign clients whose data naturally drifts from the global distribution.

First, we analyze performance under attack by observing the trend as non-IID skew increases ($\alpha$ from 1.0 to 0.1), shown in Figures~\ref{fig:non_iid_robustness1} and~\ref{fig:non_iid_robustness2}. While the undefended model (No-Def) collapses, P4P's defensive accuracy remains remarkably stable across all $\alpha$ levels. This is most evident in the most severe setting ($\alpha=0.1$), where P4P maintains high accuracy (approx. 90\%) under attacks like Sign Flip. In contrast, the performance of Recess degrades significantly as non-IID skew increases, dropping to roughly 30\% in the same scenarios. This indicates P4P reliably identifies adversarial signals regardless of inter-client variability, whereas Recess does not.

Second, we address the critical challenge of false positives by analyzing the "No Attack" scenario trend in Table~\ref{tab:fed_performance}. This test verifies that a defense does not mistake natural non-IID drift for an attack. At $\alpha=1.0$ (IID), P4P (95.85\%) and Recess (90.23\%) both perform well. However, as heterogeneity increases to $\alpha=0.5$, Recess's accuracy drops to 79.87\%, and at $\alpha=0.1$, it collapses to 59.08\%. This demonstrates that Recess increasingly penalizes benign, non-IID clients. In stark contrast, P4P's accuracy (87.55\% at $\alpha=0.1$) remains stable and nearly identical to the No-Def baseline (90.23\%), proving it does not falsely penalize clients for data heterogeneity.

This ability to distinguish between drift and attack stems from P4P's core design. Defenses like Recess, which rely on majority consensus, compare clients \textit{to each other}. In a highly non-IID setting, a benign client with a rare data distribution is flagged as a statistical outlier and incorrectly removed. P4P avoids this trap. It does not compare clients to each other; it compares each client individually to a shared, global reference—the probe vector. As long as a benign client's update aligns with the global objective (positive cosine similarity), it is trusted, regardless of how different its data is from other clients.

These findings comprehensively answer RQ2. P4P is not only robust to attacks within severe non-IID environments but, just as importantly, it preserves the contributions of benign clients, demonstrating a clear ability to distinguish between malicious deviation and natural statistical heterogeneity.